\section{Introduction}
\label{sec:introduction}

\subsection{Background and Motivation}
\label{sec:background_and_motivation}

\IEEEPARstart{R}{eaching} a net-zero economy by 2050 requires extensive electrification of the energy system. During the past two decades, renewable generation, particularly wind and photovoltaic power, has grown steadily and is expected to accelerate towards 2050~\cite{international_energy_agency_world_2023}. At the same time, decarbonization targets require the electrification of transport, heating, and industrial processes, replacing fossil fuels with electricity. In Europe, studies project that the total electricity demand will triple by 2050~\cite{energyville_paths_2023}. Consequently, future power systems will be largely dominated by power-electronic interfaced resources on both the supply and demand sides. This transition fundamentally changes the harmonic behavior of power systems. Traditional grids, dominated by synchronous machines and mostly linear loads, exhibited limited harmonic distortion. Modern systems, however, increasingly rely on converters whose nonlinear characteristics generate a broad spectrum of harmonic currents. The aggregation of numerous distributed converters creates new challenges in maintaining harmonic voltage quality across transmission and distribution networks.

A key challenge arises from the stochastic nature of harmonic emissions. Measurements and laboratory studies show that harmonic magnitudes depend on loading, switching states, and internal control actions. More importantly, harmonic current phase angles, which govern the phasor summation of emissions from multiple sources, vary significantly. Even small angle variations can result in significant differences in harmonic voltage distortion, particularly in networks with many distributed converters. As the penetration of the converter increases, system operators need transparent metrics to assess whether additional units can be connected without exceeding harmonic limits. Harmonic hosting capacity provides such a metric by defining the maximum allowable harmonic emissions while ensuring compliance with power-quality standards. However, most existing methods rely on deterministic assumptions regarding harmonic magnitudes, phase angles, network parameters, and background distortion, leading to overly conservative or insufficiently protective results.

\subsection{Literature Review}
\label{sec:literature_review}

A growing body of literature presents approaches to determine the harmonic hosting capacity of a power system and to characterize uncertainty in harmonic power flow. Table~\ref{tab:literature_review} summarizes the publications discussed below, characterized by their mathematical method, network representation, stochastic method and tractability, and by which of the following features they represent stochastically:
\begin{enumerate*}
    \item [(a)] unit harmonic injection,
    \item [(b)] background harmonic distortion,
    \item [(c)] network parameters,
    \item [(d)] network exploitation, and
    \item [(e)] harmonic limits.
\end{enumerate*}

\begin{table*}[!t]
    \centering
    \caption{Literature review on harmonic hosting capacity and stochastic harmonic power flow, in order of discussion, enumerating the stochastic features~(a)--(e). The entry marked $\star$ is the deterministic predecessor of this work.}
    \label{tab:literature_review}
    \begin{tabular}{ c | c | c | c c | c | c c c c c | r r }
        \multirow{2}{*}{reference} & \multirow{2}{*}{objective} & \multirow{2}{*}{method$^{a}$} & \multicolumn{2}{c |}{network} & \multirow{2}{*}{stochastic$^{b}$} & \multicolumn{5}{c |}{stochastic features} & \multicolumn{2}{c}{tractability$^{c}$} \\
                                            &       &       & model & phases &       & (a)  & (b)   & (c)   & (d)   & (e)   & nodes & minutes \\ \hline
        \cite{noauthor_iec61000-3-6_2008}   & hhc   & ae    & nodal & bal.  & an.   &\cmark &\xmark &\xmark &\xmark &\cmark & --    & --    \\
        \cite{santos_considerations_2015}   & hhc   & ae    & nodal & bal.  & an.   &\cmark &\cmark &\xmark &\xmark &\xmark & --    & --    \\
        \cite{cho_analysis_2019}            & hhc   & pf    & full  & bal.  & mc    &\cmark &\xmark &\cmark &\xmark &\xmark & --    & --    \\ \arrayrulecolor{gray!20} \hline
        \cite{li2020}                       & hpf   & pf    & full  & bal.  & mc    &\cmark &\xmark &\cmark &\xmark &\xmark & 14    & --    \\
        \cite{xie_piecewise_2022}           & hpf   & pf    & full  & unb.  & an.   &\cmark &\xmark &\xmark &\xmark &\xmark & 8500  & 25    \\
        \cite{xie2022}                      & hpf   & pf    & full  & unb.  & an.   &\cmark &\xmark &\xmark &\xmark &\xmark & 8500  & 119   \\
        \cite{cadenazarate2024}             & hpf   & pf    & full  & bal.  & mc    &\cmark &\xmark &\xmark &\xmark &\xmark & --    & --    \\ \arrayrulecolor{gray!20} \hline
        \cite{sun_distribution_2012}        & hhc   & nl    & full  & bal.  & --    &\xmark &\xmark &\xmark &\xmark &\xmark & 5     & --    \\
        \cite{sun_incorporating_2012}       & hhc   & nl    & full  & bal.  & --    &\xmark &\xmark &\xmark &\xmark &\xmark & --    & --    \\
        \cite{sakar_integration_2018}       & hhc   & pf    & full  & bal.  & --    &\xmark &\xmark &\xmark &\xmark &\xmark & 3     & --    \\
        \cite{sun_distribution_2022}        & hhc   & cc    & full  & bal.  & scen. &\xmark &\cmark &\xmark &\cmark &\cmark & 16    & 30    \\
        \cite{nakhodchi2022}                & hhc   & ae    & nodal & bal.  & mc    &\cmark &\xmark &\xmark &\xmark &\cmark & --    & --    \\
        \cite{nakhodchi_impact_2023}        & hhc   & ae    & nodal & bal.  & mc    &\cmark &\xmark &\xmark &\xmark &\xmark & 64    & --    \\ \arrayrulecolor{black} \hline
        $\star$~\cite{vanacker_hosting_2026}& hhc   & soc   & full  & bal.  & --    &\xmark &\xmark &\xmark &\xmark &\xmark & 1880  & $<$\,6 \\ \arrayrulecolor{black} \hline
        This work                           & hhc   & qcqp/soc & full  & bal.  & pce   &\cmark &\xmark &\xmark &\xmark &\cmark & 30    & $<$\,0.04 \\
    \end{tabular}

    \vspace{2pt}
    \parbox{\textwidth}{\scriptsize
    $^{a}$Mathematical methods: (ae)~analytical expression, (pf)~power flow, (nl)~nonlinear optimal power flow, (soc)~second-order cone optimal power flow, (cc)~chance-constrained optimal power flow, (qcqp)~quadratically constrained quadratic program. Objective: (hhc)~harmonic hosting capacity, (hpf)~harmonic power flow. Phases: (bal.)~balanced, (unb.)~unbalanced.
    $^{b}$Stochastic methods: (mc)~Monte Carlo simulation, (an.)~analytical or moment-based, (scen.)~scenario-based, (pce)~polynomial chaos expansion.
    $^{c}$Tractability expressed as the number of nodes and the solve time in minutes.}
\end{table*}

\subsubsection{Standards and the summation law}

%The standard IEC~61000-3-6 allocates planning levels over installations connected to a point of common coupling and aggregates their emissions with a general summation law, $I_{\Harmonic} = (\sum_{\Unit} I_{\Unit,\Harmonic}^{\alpha})^{1/\alpha}$, with indicative exponents $\alpha = 1$ for $\Harmonic < 5$, $1.4$ for $5 \leq \Harmonic \leq 10$ and $2$ for $\Harmonic > 10$, and compares the 95\,\% distortion values with the planning levels~\cite{noauthor_iec61000-3-6_2008}. The summation law thus carries the diversity of the harmonic phase angles into the standard, but its exponents are indicative. Santos et al.\ used the summation law to determine the harmonic hosting capacity of transmission and distribution systems, represented the upstream network by a Norton equivalent with background distortion of unknown phase, and obtained a harmonic hosting capacity between 37\,MW and 160\,MW for the in-phase and anti-phase extremes; they also cautioned that the indicative exponents may under- or overestimate the aggregate of correlated sources such as wind and photovoltaic plants~\cite{santos_considerations_2015}. Field measurements confirm this caution: Wang et al.\ analyzed synchronized measurements in North-American residential distribution systems and found that low-order harmonic currents from different homes tend to add up, i.e., the effects is much less than claimed in the literature, while recombination occurs at high orders only~\cite{wang_characteristics_2017}. Cho et al.\ compared IEC~61000-3-6 with the deterministic current limits of IEEE Standard~519, set as a function of the short-circuit ratio~\cite{noauthor_ieee519_2022}, by Monte Carlo simulation on a multi-feeder distribution model, and showed that the latter can drive the voltage distortion of distribution systems up to twice the planning levels~\cite{cho_analysis_2019}. Trelles-Molina et al.\ extended the summation law from transmission systems to individual customer facilities with few harmonic sources and analyzed the probabilities that support the exponents~\cite{trellesmolina_extension_2022}. How much recombination a given network and converter population actually allow therefore remains an open question.
The standard IEC~61000-3-6 allocates planning levels over installations connected to a point of common coupling and aggregates their emissions, i.e., combines the harmonic currents of the individual installations into a single total current at the point of common coupling, with a general summation law, $I_{\Harmonic} = (\sum_{\Unit} I_{\Unit,\Harmonic}^{\alpha})^{1/\alpha}$, with indicative exponents $\alpha = 1$ for $\Harmonic < 5$, $1.4$ for $5 \leq \Harmonic \leq 10$ and $2$ for $\Harmonic > 10$, and compares the 95\,\% distortion values with the planning levels~\cite{noauthor_iec61000-3-6_2008}. The summation law thus carries the diversity of the harmonic phase angles into the standard, but its exponents are indicative. Santos et al.\ used the summation law to determine the harmonic hosting capacity of transmission and distribution systems, represented the upstream network by a Norton equivalent with background distortion of unknown phase, and obtained a harmonic hosting capacity between 37\,MW and 160\,MW for the in-phase and anti-phase extremes; they also cautioned that the indicative exponents may under- or overestimate the aggregate of correlated sources such as wind and photovoltaic plants~\cite{santos_considerations_2015}. Field measurements confirm this caution: Wang et al.\ analyzed synchronized measurements in North-American residential distribution systems and found that low-order harmonic currents from different homes tend to add up, i.e., the effects is much less than claimed in the literature, while recombination occurs at high orders only~\cite{wang_characteristics_2017}. Cho et al.\ compared IEC~61000-3-6 with the deterministic current limits of IEEE Standard~519, set as a function of the short-circuit ratio~\cite{noauthor_ieee519_2022}, by Monte Carlo simulation on a multi-feeder distribution model, and showed that the latter can drive the voltage distortion of distribution systems up to twice the planning levels~\cite{cho_analysis_2019}. Trelles-Molina et al.\ extended the summation law from transmission systems to individual customer facilities with few harmonic sources and analyzed the probabilities that support the exponents~\cite{trellesmolina_extension_2022}. How much recombination a given network and converter population actually allow therefore remains an open question.

%\subsubsection{Stochastic harmonic power flow}
\subsubsection{Probabilistic harmonic power flow}

Probabilistic harmonic power flow quantifies the uncertainty of the harmonic emissions and the associated network response. Li et al.\ separated uncertainty sources due to natural variability, i.e., aleatory, from those due to a lack of knowledge, i.e., epistemic, defined quantitative amplification indices as ratios of transfer to driving-point impedances, and assessed them with Monte Carlo simulation~\cite{li2020}. Monte Carlo and data-driven methods have since become the standard tool for probabilistic harmonic power flow: Xie and Sun proposed a piecewise probabilistic harmonic power flow for unbalanced residential distribution systems~\cite{xie_piecewise_2022}, Xie et al.\ a data-driven approach for distribution systems with photovoltaic generation~\cite{xie2022}, and Cadena-Zarate et al.\ a simulation-based approach for low-voltage networks with distributed energy resources~\cite{cadenazarate2024}. These methods quantify uncertainty, but do not optimize the hosting capacity.

\subsubsection{Stochastic harmonic hosting capacity}

Early harmonic hosting capacity methods incorporated harmonic limits within deterministic frameworks. Sun et al.\ incorporated harmonic voltage constraints into an optimal power flow for distribution network capacity assessment~\cite{sun_distribution_2012} and extended it to active networks~\cite{sun_incorporating_2012}, and Sakar et al.\ analyzed the hosting capacity of large photovoltaic plants in the presence of background distortion~\cite{sakar_integration_2018}. Sun and Harrison proposed the first optimization-based stochastic harmonic hosting capacity method, with percentile-based harmonic distortion limits as chance constraints within an ac optimal power flow, and showed that relaxing the risk level increases the admissible hosting capacity~\cite{sun_distribution_2022}. Nakhodchi et al.\ included the measured emission uncertainty of a fast charging station for electric vehicles through Bayesian statistics and the correlation between harmonic orders~\cite{nakhodchi2022}, and showed that the estimated hosting capacity depends strongly on how the medium-voltage network and the remote loads are modeled~\cite{nakhodchi_impact_2023}. The present authors determined an exact lower bound for the equitable harmonic hosting capacity with a second-order cone formulation of the full network, but with fixed harmonic current angles~\cite{vanacker_hosting_2026}, and studied the impact of underground cables on it~\cite{vanacker_impact_2025}.

\subsubsection{Research gaps}

Table~\ref{tab:literature_review} shows that no existing method simultaneously achieves:
\begin{enumerate*}
%    \item [(i)] tractable network-wide optimization,
    \item [(i)] computationally tractable network-wide optimization,
    \item [(ii)] explicit treatment of phase angle uncertainty, and
    \item [(iii)] chance-constraint enforcement of harmonic distortion limits.
\end{enumerate*}
Moreover, the recombination that a network and a converter population allow has not been quantified against the summation law of IEC~61000-3-6, and the conditions under which it occurs have not been identified.

\subsection{Contributions and Outline}
\label{sec:contribution_and_outline}

%In this paper, a holistic one-shot methodology is presented to determine the steady-state harmonic hosting capacity for an entire phase-balanced three-phase power system in a stochastic context. The presented method considers uncertainty on the angle of the injected harmonic current, enabling harmonic recombination, and limits the resulting individual and total harmonic voltage distortion, as well as the root-mean-square voltages and currents, using chance constraints. The problem is formulated from the perspective of the power system operator, who may anticipate the converter technologies its customers will connect, but not the operating point, parameters and control of the individual units. The angle of the harmonic current is therefore modeled as a random variable per unit, characterized for the anticipated technology, whereas its magnitude is the decision, and the harmonic hosting capacity is determined in terms of harmonic current magnitude using a fairness principle. The presented formulation uses the rectangular current-voltage variable space and incorporates:
In this paper, a holistic single-run methodology is presented to determine the steady-state harmonic hosting capacity for an entire phase-balanced three-phase power system in a stochastic context. The presented method considers uncertainty on the angle of the injected harmonic current, enabling harmonic recombination, and limits the resulting individual and total harmonic voltage distortion, as well as the root-mean-square voltages and currents, using chance constraints. The problem is formulated from the perspective of the power system operator, who may anticipate the converter technologies its customers will connect, but not the operating point, parameters and control of the individual units. The angle of the harmonic current is therefore modeled as a random variable per harmonic source unit, characterized for the anticipated technology, whereas its magnitude is the decision: the harmonic hosting capacity is determined as the harmonic current magnitude each unit may inject, using a fairness principle. The presented formulation uses the rectangular current-voltage variable space and incorporates:
\begin{itemize}
%    \item uncertainty on the angle of the injected harmonic current via polynomial chaos expansion,
    \item uncertainty on the angle of the injected harmonic current by applying the polynomial chaos expansion technique~\cite{sullivan_introduction_2015},
    \item exact representation of the harmonic impedance in the frequency domain and detailed treatment of transformers,
    \item enforcement of stochastic harmonic limits on bus voltage and component ampacity through chance constraints, and
    \item equitable distribution of the harmonic hosting capacity following a fairness principle.
\end{itemize}
%For fairness criteria that fix the ratios between the allocations of the units, e.g., absolute equality, the problem reduces to a second-order cone program with one pair of variables per harmonic. %
For fairness criteria that fix the ratios between the allocations of the units, e.g., absolute equality, which grants every unit the same harmonic current~\cite{vanacker_hosting_2026}, the problem reduces to a second-order cone program with one pair of variables per harmonic, i.e., a convex problem that standard solvers handle efficiently, as detailed in Section~\ref{sec:absolute_equality_socp}. %
The models and numerical illustrations are made freely available as part of the \textsc{Julia}-package \textsc{HarmonicPowerModels.jl}.\footnote{Available at \url{https://github.com/timmyfaraday/HarmonicPowerModels.jl}.}

The remainder of the paper is organized as follows: Section~\ref{sec:uncertainty} addresses the uncertainty of harmonic current phasors associated with power-electronic interfaced resources. Section~\ref{sec:mathematical_framework} presents the mathematical formulation of the stochastic harmonic hosting capacity problem. Section~\ref{sec:numerical_illustration} proves the equivalence of the presented methods on the IEEE 30-bus system and investigates the levers for harmonic recombination. Section~\ref{sec:conclusion} concludes the paper.
